\chapter{Introducción}
\label{chapter:introduccion}


%%% SECTION
\section{Descripción general del problema}

En la actualidad, los procesos de minería de datos requieren grandes cantidades de datos, que en muchas ocasiones contienen información personal y privada de usuarios o personas. Aunque se realicen procesos básicos de anonimización sobre los datos, es decir, eliminación de los nombres u otros identificadores clave, existen multitud de técnicas de re-identificación que permiten volver a identificar a un usuario dentro de este conjunto de datos. En la Figura \ref{fig:context-anoni1} se presenta un mapa donde es posible contextualizar los procesos de anonimización y re-identificación dentro de un proceso de minería de datos.

\begin{figure}
	\centering
	\includegraphics[width=0.6\textwidth]{figs/image1.png}
	\caption{Pie de la imagen.}
	\label{fig:context-anoni1}
\end{figure}

\subsection{Ejemplo de subsection}

Aunque se han realizado importantes avances en preservación de la privacidad en publicación de datos, tales como el modelo \textit{k}-anonymity \cite{Sweeney:2002}.

Un ejemplo de pseudo-código se puede encontrar en el Código \ref{code:RandomSwitch-1}

\begin{algorithm}
	\caption{Pseudocódigo del algoritmo \textit{Random Switch}}
	\label{code:RandomSwitch-1}
	\begin{algorithmic}
		\REQUIRE{El grafo original $G$ y el porcentaje de anonimización $p$ que se desea aplicar.}
		\ENSURE{El grafo $G$ anonimizado.}
		\STATE $num = round(G.num\_edges() * p)$
		\STATE $i = 0$
		\WHILE {$i < num$}
		\STATE {$e_{1} = G.random\_edge()$}
		\STATE $e_{2} = G.random\_edge()$
		\STATE $new\_e_{1} = (e_{1}.origen, e_{2}.origen)$
		\STATE $new\_e_{2} = (e_{1}.destino, e_{2}.destino)$
		\IF {$!G.exist(new\_e_{1})$ \AND $!G.exist(new\_e_{2})$}
		\STATE $G.add\_edge(new\_e_{1})$
		\STATE $G.add\_edge(new\_e_{2})$
		\STATE $G.delete\_edge(e_{1})$
		\STATE $G.delete\_edge(e_{2})$
		\STATE $i=i+1$
		\ENDIF
		\ENDWHILE
		\RETURN $G$
	\end{algorithmic}
\end{algorithm}

Un ejemplo de tabla se puede ver en la Tabla \ref{table:ejemplo_vertex_refi_query}

\begin{table}
	\centering{}
	\begin{tabular}{ l || c | c | l }
		\hline
		Node ID & $\mathcal{H}_{0}$ & $\mathcal{H}_{1}$ & $\mathcal{H}_{2}$ \\
		\hline
		\hline
		Alice & $\epsilon$ & 1 & \{4\}  \\
		\hline
		Bob & $\epsilon$ & 4 & \{1, 1, 4, 4\}  \\
		\hline
		Carol & $\epsilon$ & 1 & \{4\}  \\
		\hline
		Dave & $\epsilon$ & 4 & \{2, 4, 4, 4\}  \\
		\hline
		Ed & $\epsilon$ & 4 & \{2, 4, 4, 4\}  \\
		\hline
		Fred & $\epsilon$ & 2 & \{4, 4\}  \\
		\hline
		Greg & $\epsilon$ & 4 & \{2, 2, 4, 4\}  \\
		\hline
		Harry & $\epsilon$ & 2 & \{4, 4\}  \\
		\hline
	\end{tabular}
	\caption{\textit{Vertex refinement queries}.}
	\label{table:ejemplo_vertex_refi_query}
\end{table}





\section{Introducción v2}

Esta plantilla pretende ser una guía para los estudiantes. Esta plantilla se puede adaptar a las necesidades específicas de cada proyecto si el supervisor del proyecto está de acuerdo con los cambios.

\begin{figure}[h]
\centering
\includegraphics[width=0.5\textwidth]{./figs/image1.png}
\caption{Ejemplo de figura. Será indexada en la “Lista de Figuras”.}
\label{fig:figura_ejemplo}
\end{figure}

\subsection{Contexto y motivación}

Punto de partida del proyecto (¿Cuál es el problema que necesita ser resuelto? ¿Por qué es un tema relevante? ¿Cómo se está resolviendo el problema actualmente?) y descripción de la contribución (¿Qué resultado se espera obtener?)

\subsection{Objetivos}

Listado de los objetivos a alcanzar en este proyecto.

\subsection{Sostenibilidad, diversidad y desafíos ético/sociales}

Esta sección debe evaluar el impacto positivo/negativo del proyecto en las siguientes dimensiones. No es necesario alcanzar un impacto positivo en todas las dimensiones, pero es necesario considerar y discutir si existe un impacto desde el inicio del proyecto.

\begin{description}
    \item[Sostenibilidad] En el desarrollo del proyecto o durante todo su ciclo de vida (por ejemplo, despliegue, retiro), ¿tiene el resultado de este proyecto un impacto en la sostenibilidad y/o huella ecológica (consumo/ahorro de energía/recursos, desperdicio, contaminación, agotamiento de materias primas)? ¿Está afectado por leyes o regulaciones sobre este asunto? Considerando otra perspectiva, ¿afecta a alguno de los Objetivos de Desarrollo Sostenible (ODS) relacionados con estas dimensiones? Si no tiene ningún impacto, ya sea positivo o negativo, debe explicar cómo llegó a esta conclusión y justificar su respuesta.
    \item[Comportamiento ético y responsabilidad social] ¿Es el resultado del proyecto demasiado técnico para tener algún impacto positivo/negativo en aspectos éticos/sociales? ¿Tiene un impacto en leyes/regulaciones (datos, privacidad, trabajo, propiedad intelectual, seguridad personal, …)? ¿Se adhiere a los principios deontológicos de la profesión? ¿Pone en peligro/mejora/empeora algún puesto de trabajo? Si no tiene ningún impacto, ya sea positivo o negativo, debe explicar cómo llegó a esta conclusión y justificar su respuesta.
    \item[Diversidad, género y derechos humanos] ¿Es el resultado de este proyecto tan técnico que no tiene impacto positivo/negativo en términos de género, diversidad o derechos humanos? ¿Y en leyes/regulaciones? ¿Y en términos de accesibilidad, discapacidad, ergonomía y/o seguridad de la información? Si no tiene ningún impacto, ya sea positivo o negativo, debe explicar cómo llegó a esta conclusión y justificar su respuesta.
\end{description}

\subsection{Enfoque y metodología}

Describa las estrategias potenciales para desarrollar este proyecto y explique la estrategia seleccionada. Discuta por qué la estrategia seleccionada es la más adecuada para alcanzar los objetivos del proyecto.

\subsection{Planificación}

Descripción de los recursos necesarios para desarrollar el proyecto, las tareas a realizar y una planificación temporal de cada tarea utilizando un diagrama de Gantt o equivalente. Esta planificación debe definir los hitos que se completarán en cada Prueba de Evaluación Continua (PEC).

\subsection{Resumen de los productos del proyecto}

No es necesario describir cada producto en detalle: esto se hará en los capítulos restantes del proyecto.

\subsection{Breve descripción de los demás capítulos del informe}

Breve descripción de los contenidos de cada capítulo y su relación con el resto del proyecto.
